\section{DeepSeek V3 y R1: análisis técnico}

\subsection{Panorama del modelo}

DeepSeek V3 es un modelo de lenguaje Transformer de Mixture-of-Experts (MoE) de la empresa china DeepSeek, de pesos abiertos y en su versión con fine-tuning por instrucciones. R1 es un modelo de razonamiento entrenado a partir de V3 mediante aprendizaje por refuerzo, publicado el 20 de enero de 2,025. Ambos comparten el mismo modelo base preentrenado, pero sus rutas de posentrenamiento difieren.

\textit{«DeepSeek ha hecho un trabajo excelente para difundir la comprensión de la IA: sus artículos son extremadamente detallados.»} (Nathan) Este nivel de detalle es sumamente inusual entre los laboratorios de vanguardia: los artículos incluyen curvas de loss completas del entrenamiento, justificaciones de la elección de hiperparámetros y experimentos de ablación de arquitectura, lo que permite que otros equipos los retomen directamente.

Contexto: la empresa matriz de DeepSeek es Huanfang Lianghua (幻方量化, Highflyer), un fondo de cobertura cuantitativo chino cuyo CEO, Liang Wenfeng (梁文锋), dirige ambas empresas a la vez. Dylan lo compara con \textit{«una figura al estilo Elon/Jensen—que participa en todo, con un carácter totalmente AGI»}. Highflyer acumuló 10,000 GPU A100 antes de los controles de exportación de 2,021, una previsión que proporcionó la base de hardware clave para el DeepSeek posterior.

\begin{importantbox}{Diferencias en la experiencia de usuario: V3 frente a R1}
\textbf{V3}: modelo de chat estándar, de generación rápida y salida en Markdown de alta calidad.

\textbf{R1}: muestra un proceso extendido de razonamiento Chain-of-Thought—el modelo descompone el problema, reflexiona sobre sí mismo, retrocede para corregirse y después da la respuesta. Este proceso de razonamiento visible cautivó la imaginación del público.

Ejemplo: al preguntarle a R1 por «una idea realmente novedosa sobre los seres humanos», el modelo razonó durante 157 segundos y finalmente respondió: \textit{«los seres humanos convierten instintivamente los deseos egoístas en sistemas de cooperación fingiendo colectivamente que reglas abstractas (el dinero, la ley, los derechos) son reales.»}

La evaluación poética de Lex sobre el Chain-of-Thought de R1: \textit{«su no linealidad se parece al Ulises de James Joyce—salta entre el chino y el inglés, produce fragmentos que parecen ser ruido y luego, de pronto, da una respuesta clara.»}

En las pruebas comparativas, O1 Pro fue sistemáticamente el mejor en preguntas filosóficas, R1 quedó en segundo lugar y Gemini Flash en tercero. Pero el proceso de razonamiento visible de R1 crea una experiencia estética singular. Gemini Flash Thinking, de Google, probablemente sea más barato que R1 y no menos capaz, pero casi nadie habla de él.
\end{importantbox}

\begin{figure}[H]
\centering
\includegraphics[width=0.92\textwidth]{frame-01-open-weights.jpg}
\caption{Fragmento de discusión sobre el Chain-of-Thought visible de R1\protect\footnotemark}
\end{figure}
\footnotetext{Intervalo de tiempo del video: 00:12:00--00:12:10.}

\subsection{Pesos abiertos frente a código abierto}

\begin{knowledgebox}{El espectro de la apertura}
\begin{itemize}
    \item \textbf{Pesos abiertos}: los pesos del modelo se pueden descargar y ejecutar localmente—tus datos permanecen en tu equipo
    \item \textbf{Código completamente abierto}: pesos + datos de entrenamiento + código de entrenamiento (el estándar de AI2)
    \item DeepSeek R1: licencia MIT—sin restricciones comerciales posteriores y utilizable para generar datos sintéticos
    \item Llama: licencia más restrictiva (limitaciones de casos de uso, requisitos de marca)
\end{itemize}

\textit{«quien te roba los datos no es el modelo, sino quien lo aloja.»} (Nathan)

Punto clave: los artículos de DeepSeek son extremadamente detallados y ofrecen detalles de entrenamiento aplicables en la práctica (incluidas las curvas de loss), pero los datos y el código de entrenamiento no se hicieron públicos. Esto es «pesos abiertos» y no «código completamente abierto»—pero aun así, es el mayor grado de apertura entre los modelos de vanguardia.

La licencia MIT de DeepSeek R1 es un «gran reinicio»—el primer modelo de vanguardia con una licencia verdaderamente permisiva. Antes de esto, las opciones disponibles o no eran de vanguardia, o tenían licencias restrictivas (como los requisitos de marca y las limitaciones de uso de Llama). \textit{«Necesitamos modelos verdaderamente abiertos… puedes hacer con R1 una copia barata y hacer como si fuera tuya. Llama no permite eso.»} (Nathan)
\end{knowledgebox}

\subsection{El secreto del bajo costo de entrenamiento}

DeepSeek afirma haber entrenado V3 con 2,000 GPU H800 (solo la etapa de preentrenamiento). Dos innovaciones clave de arquitectura:

\begin{enumerate}
    \item \textbf{MoE (Mixture of Experts)}: 671B parámetros en total, pero cada token solo activa alrededor de 37B (8 de los 256 expertos)—lo que reduce drásticamente el cómputo. Esta no es una técnica inventada por DeepSeek (el Switch Transformer de Google es anterior), pero DeepSeek la llevó al extremo en un modelo de vanguardia.
    \item \textbf{MLA (Multi-head Latent Attention)}: proyecta los valores Key-Value del mecanismo de atención a un espacio latente de baja dimensión, lo que ahorra entre 80 y 90\% de memoria. Esto hace posible procesar contextos más largos dentro de la memoria limitada de las H800.
\end{enumerate}

Pero el verdadero «as bajo la manga» está en un nivel aún más bajo—DeepSeek hizo modificaciones por debajo de la capa CUDA, programando manualmente los núcleos SM (Streaming Multiprocessor) para la comunicación, evitando la biblioteca de comunicación estándar de NVIDIA, NCCL.

\textit{«la necesidad es la madre de la innovación—tuvieron que hacerlo porque les recortaron el ancho de banda de interconexión.»} (Dylan) La única diferencia entre H800 y H100 es que se redujo el ancho de banda de interconexión NVLink, y la comunicación entre nodos es justamente el cuello de botella del entrenamiento a gran escala. DeepSeek se vio obligado a innovar en este punto restringido, y terminó desarrollando estrategias de comunicación más eficientes que las soluciones estándar.

\begin{warningbox}{La verdad detrás de las «2,000 GPU»}
Las 2,000 H800 que menciona el artículo \textbf{se refieren únicamente a una ejecución del preentrenamiento de V3}. SemiAnalysis calcula que DeepSeek en realidad cuenta con alrededor de 50,000 GPU—porque además se necesitan:
\begin{itemize}
    \item la búsqueda de arquitectura previa y los experimentos de ablación (ablation studies) a pequeña escala
    \item el posentrenamiento (SFT + RLHF)
    \item el entrenamiento de RL específico para R1
    \item el despliegue del servicio de inferencia
\end{itemize}

Comparación: Meta cuenta con más de 400,000 GPU; Llama 3 se entrenó con 16,000. La verdadera ventaja de DeepSeek no es «lo barato», sino \textbf{su capacidad de optimización extrema en hardware restringido}.

Contexto: la empresa matriz de DeepSeek, Highflyer (幻方量化), es un fondo de cobertura cuantitativo chino que acumuló 10,000 GPU A100 antes de los controles de exportación de 2,021. Liang Wenfeng, su CEO, es comparado por Dylan con \textit{«una figura al estilo Elon/Jensen—que participa en todo, con un carácter totalmente AGI»}.
\end{warningbox}

\subsubsection{Bitter Lesson y la filosofía de entrenamiento}

El éxito de DeepSeek encarna el Bitter Lesson de Rich Sutton: los métodos escalables en aprendizaje y búsqueda ganan a largo plazo; hay que minimizar los sesgos previos humanos. La innovación se acumula con el tiempo: las pequeñas mejoras en los datos, la arquitectura y el posentrenamiento se suman continuamente.

El momento decisivo clave del entrenamiento es el «YOLO run»—después de las ablaciones a pequeña escala, se destinan todos los recursos a una gran ejecución de entrenamiento. El GPT-4 YOLO de OpenAI en 2,022 fue el más audaz: una arquitectura MoE totalmente nueva, con todo el cómputo dedicado durante meses.

\textit{«los modelos solo quieren aprender—hay que darles un paisaje de pérdida simple y quitarles los obstáculos del camino.»} (Nathan)

\subsubsection{Loss Spike durante el entrenamiento}

Los picos de loss durante el entrenamiento son un problema que enfrentan todos los laboratorios: algunos provienen de datos defectuosos (un ejemplo es el subreddit «pandilla del microondas»), otros de inestabilidad numérica. Los picos rápidos y los picos lentos requieren estrategias de recuperación distintas.

Los ingenieros vigilan continuamente las curvas de loss incluso durante la cena, revisando el teléfono cada 10 minutos. Todas las empresas tienen ejecuciones de entrenamiento fallidas—ese es el costo de avanzar la frontera. El entrenamiento en FP8 introduce aún más inestabilidad.

\subsection{Resumen de la sección}

El avance central de DeepSeek es la optimización de ingeniería llevada al extremo bajo las restricciones de los controles de exportación: el MoE reduce los parámetros activados, el MLA reduce la memoria y la personalización a bajo nivel de CUDA. El nivel de detalle del artículo de V3 ofrece a toda la industria una hoja de ruta técnica aplicable en la práctica.
